% chapters/chapter2/2_2_lagrange.tex
% !TeX root=../../main.tex

\section{معرفی سیستم}

در این سیستم بسط‌یافته، یک آونگ ساده به ارابه‌ای متصل است که روی یک ریل افقی حرکت می‌کند. طول کابل آونگ متغیر است و به عنوان یک متغیر سینماتیکی (ورودی) با نماد $l(t)$ در نظر گرفته می‌شود. همچنین، نیروی محرک $F$ سیستم را در امتداد ریل تحریک می‌کند (جهت مثبت این نیرو، هم‌سو با جهت مثبت محور $x$ فرض شده است). شماتیک سینماتیکی و پارامترهای مدل‌سازی سیستم در شکل \ref{fig:modeling_schematic} نشان داده شده است.

\begin{figure}[htbp]
    \centering
    \includegraphics[width=0.7\textwidth]{Modeling.pdf}
    \caption{شماتیک سینماتیکی جرثقیل سقفی با طول کابل متغیر و درجات آزادی سیستم.}
    \label{fig:modeling_schematic}
\end{figure}

مختصات تعمیم‌یافته سیستم با توجه به دو درجه آزادی آن، به صورت زیر انتخاب می‌شوند:
\[
    q = \begin{bmatrix} x \\ \theta \end{bmatrix}
\]
متغیرها و پارامترهای سیستم عبارت‌اند از:
\begin{itemize}
    \item $x$: موقعیت افقی ارابه نسبت به مبدأ مختصات
    \item $\theta$: زاویه آونگ نسبت به راستای قائم (جهت پادساعت‌گرد، مثبت فرض شده است)
    \item $l(t)$: طول متغیر کابل آونگ (ورودی سینماتیکی)
    \item $M$: جرم ارابه
    \item $m$: جرم نقطه‌ای آونگ
    \item $g$: شتاب گرانش
    \item $b_c$: ضریب اصطکاک لزج میان ارابه و ریل
    \item $b_p$: ضریب اصطکاک لزج در مفصل آونگ
\end{itemize}

\section{استخراج معادلات حرکت به روش لاگرانژ}

\subsection{سینماتیک و مختصات جرم آونگ}

با در نظر گرفتن دستگاه مختصات دکارتی استاندارد (محور $x$ به سمت راست و محور $y$ به سمت بالا)، مختصات جرم آونگ به صورت زیر بیان می‌شود:
\begin{align*}
    x_p & = x + l(t) \sin\theta \\
    y_p & = -l(t) \cos\theta
\end{align*}
با مشتق‌گیری زمانی از روابط فوق (با فرض $\dot{l} \neq 0$)، مؤلفه‌های سرعت محاسبه می‌شوند:
\begin{align*}
    \dot{x}_p & = \dot{x} + \dot{l}\sin\theta + l\dot{\theta}\cos\theta \\
    \dot{y}_p & = -\dot{l}\cos\theta + l\dot{\theta}\sin\theta
\end{align*}
مجذور سرعت خطی جرم آونگ ($v_p^2 = \dot{x}_p^2 + \dot{y}_p^2$) پس از بسط و ساده‌سازی عبارات، به شکل زیر به دست می‌آید:
\[
    v_p^2 = \dot{x}^2 + \dot{l}^2 + l^2\dot{\theta}^2 + 2\dot{x}\dot{l}\sin\theta + 2\dot{x}l\dot{\theta}\cos\theta
\]

\subsection{انرژی جنبشی و پتانسیل}

انرژی جنبشی کل سیستم برابر با مجموع انرژی جنبشی ارابه و آونگ است:
\begin{align*}
    T & = \frac{1}{2}M\dot{x}^2 + \frac{1}{2}m v_p^2                                                                                                          \\
    T & = \frac{1}{2}(M+m)\dot{x}^2 + \frac{1}{2}m \left( \dot{l}^2 + l^2\dot{\theta}^2 \right) + m\dot{x}\dot{l}\sin\theta + ml\dot{x}\dot{\theta}\cos\theta
\end{align*}
با فرض قرارگیری مبدأ انرژی پتانسیل در راستای حرکت ارابه، انرژی پتانسیل سیستم برابر است با:
\[
    V = -mgl\cos\theta
\]

\subsection{تابع لاگرانژ}

با استفاده از رابطه $\mathcal{L} = T - V$، تابع لاگرانژ سیستم به این شکل تشکیل می‌شود:
\[
    \mathcal{L} = \frac{1}{2}(M+m)\dot{x}^2 + \frac{1}{2}m(\dot{l}^2 + l^2\dot{\theta}^2) + m\dot{x}\dot{l}\sin\theta + ml\dot{x}\dot{\theta}\cos\theta + mgl\cos\theta
\]

\subsection{نیروهای تعمیم‌یافته غیرمحافظه‌کار}
نیروهای تعمیم‌یافته متناظر با هر درجه آزادی به صورت زیر تعریف می‌شوند:
\begin{align*}
    Q_x      & = F - b_c\dot{x}   \\
    Q_\theta & = -b_p\dot{\theta}
\end{align*}

\subsection{اعمال معادلات اویلر-لاگرانژ}

\subsubsection{معادله حرکت ارابه (مختصه $x$)}
\[
    \frac{d}{dt}\left(\frac{\partial \mathcal{L}}{\partial \dot{x}}\right) - \frac{\partial \mathcal{L}}{\partial x} = Q_x
\]
با محاسبه مشتقات جزئی، جایگذاری آن‌ها در رابطه پیشین و مرتب‌سازی عبارات، معادله حرکت ارابه به دست می‌آید:
\[
    (M+m)\ddot{x} + ml\ddot{\theta}\cos\theta + m\ddot{l}\sin\theta + 2m\dot{l}\dot{\theta}\cos\theta - ml\dot{\theta}^2\sin\theta + b_c\dot{x} = F
\]

\subsubsection{معادله حرکت آونگ (مختصه $\theta$)}
\[
    \frac{d}{dt}\left(\frac{\partial \mathcal{L}}{\partial \dot{\theta}}\right) - \frac{\partial \mathcal{L}}{\partial \theta} = Q_\theta
\]
پس از محاسبه مشتقات و ساده‌سازی عبارات، معادله دوم حاصل می‌شود:
\[
    ml^2\ddot{\theta} + ml\ddot{x}\cos\theta + 2ml\dot{l}\dot{\theta} + mgl\sin\theta + b_p\dot{\theta} = 0
\]